\documentclass{article}
\usepackage[utf8]{inputenc}
\usepackage{graphicx}
\usepackage{amsmath}
\usepackage{amssymb}
\usepackage{float} % Add the float package
\usepackage{listings}
\usepackage{xcolor}
\usepackage{array}

\title{Mathematical Statistics Assignment 4}
\author{ABID ALI (Student ID: 2023280099)}
\date{November 23, 2024}

\begin{document}

\maketitle

\newpage % Forces the following content to appear on a new page

\section*{4.1.3}

Let the population be given by $\{1, 2, 3\}$. Let $P(x) = \frac{1}{3}$ for $x = 1, 2, 3$. Take samples of size 3 with replacement.

\begin{enumerate}
    \item[(a)] Calculate $\mu$ and $\sigma^2$.
    \item[(b)] Obtain the sampling distribution of the sample mean.
    \item[(c)] Obtain the mean and variance of the sample mean.
\end{enumerate}

\section*{Solution:}


\subsection*{(a) Calculate $\mu$ and $\sigma^2$}

The population is $\{1, 2, 3\}$ with $P(x) = \frac{1}{3}$ for $x = 1, 2, 3$.

\textbf{Step 1: Calculate the population mean ($\mu$)}  
The population mean is given by:
\[
\mu = \frac{1}{N} \sum_{i=1}^N x_i
\]
where $x_i$ are the population values and $N = 3$. Substituting the values:
\[
\mu = \frac{1}{3} (1 + 2 + 3) = \frac{6}{3} = 2
\]

\textbf{Step 2: Calculate the population variance ($\sigma^2$)}  
The population variance is:
\[
\sigma^2 = \frac{1}{N} \sum_{i=1}^N (x_i - \mu)^2
\]
Substitute $\mu = 2$ and $x_i = \{1, 2, 3\}$:
\[
\sigma^2 = \frac{1}{3} \left[(1 - 2)^2 + (2 - 2)^2 + (3 - 2)^2 \right]
\]
\[
\sigma^2 = \frac{1}{3} \left[1^2 + 0^2 + 1^2\right] = \frac{1}{3} (1 + 0 + 1) = \frac{2}{3}
\]

Thus, 
\[
\mu = 2, \quad \sigma^2 = \frac{2}{3}.
\]

\subsection*{(b) Sampling distribution of the sample mean}

\textbf{Step 1: Generate all possible samples of size 3}  
Since sampling is with replacement, the total number of samples is $3^3 = 27$. The samples are:
\[
\{(1,1,1), (1,1,2), (1,1,3), (1,2,1), \dots, (3,3,3)\}.
\]

\textbf{Step 2: Calculate the sample mean for each sample}  
The sample mean for a sample $(x_1, x_2, x_3)$ is:
\[
\bar{x} = \frac{x_1 + x_2 + x_3}{3}.
\]

For example:
- For $(1, 1, 1)$, $\bar{x} = \frac{1 + 1 + 1}{3} = 1$,
- For $(1, 1, 2)$, $\bar{x} = \frac{1 + 1 + 2}{3} = 1.33$,
- For $(1, 1, 3)$, $\bar{x} = \frac{1 + 1 + 3}{3} = 1.67$.

Continue similarly for all 27 samples.

\textbf{Step 3: Tabulate the sampling distribution}  
Count the frequency of each unique sample mean and calculate their probabilities.

\[
\begin{array}{|c|c|c|c|c|c|c|}
\hline
\textbf{Sample} & \bar{X} & \textbf{Sample} & \bar{X} & \textbf{Sample} & \bar{X} \\ \hline
(1,1,1) & 1 & (2,1,1) & 1\frac{1}{3} & (3,1,1) & 1\frac{1}{2} \\ \hline
(1,1,2) & 1\frac{1}{3} & (2,1,2) & 1\frac{2}{3} & (3,1,2) & 2 \\ \hline
(1,1,3) & 1\frac{2}{3} & (2,1,3) & 2 & (3,2,1) & 2 \\ \hline
(1,2,1) & 1\frac{1}{3} & (2,2,1) & 2 & (3,2,2) & 2\frac{1}{3} \\ \hline
(1,2,2) & 2 & (2,2,2) & 2 & (3,2,3) & 2\frac{2}{3} \\ \hline
(1,2,3) & 2\frac{1}{3} & (2,3,1) & 2\frac{1}{3} & (3,3,1) & 2\frac{1}{3} \\ \hline
(1,3,1) & 2 & (2,3,2) & 2\frac{2}{3} & (3,3,2) & 2\frac{2}{3} \\ \hline
(1,3,3) & 2\frac{2}{3} & (2,3,3) & 3 & (3,3,3) & 3 \\ \hline
\end{array}
\]

\[
\begin{array}{|c|c|c|c|c|c|c|c|}
\hline
\bar{X} & 1 & 1\frac{1}{3} & 1\frac{2}{3} & 2 & 2\frac{1}{3} & 2\frac{2}{3} & 3 \\ \hline
p(\bar{X}) & \frac{1}{27} & \frac{1}{9} & \frac{2}{9} & \frac{7}{27} & \frac{2}{9} & \frac{1}{9} & \frac{1}{27} \\ \hline
\end{array}
\]


\subsection*{(c) Mean and variance of the sample mean}

\textbf{Step 1: Mean of the sample mean}  
The mean of the sampling distribution of the sample mean is:
\[
E(\bar{X}) = \mu = 2.
\]

\textbf{Step 2: Variance of the sample mean}  
The variance of the sample mean is:
\[
\text{Var}(\bar{X}) = \frac{\sigma^2}{n}.
\]
Substitute $\sigma^2 = \frac{2}{3}$ and $n = 3$:
\[
\text{Var}(\bar{X}) = \frac{\frac{2}{3}}{3} = \frac{2}{9}.
\]

Thus, 
\[
E(\bar{X}) = 2, \quad \text{Var}(\bar{X}) = \frac{2}{9}.
\]

\section*{4.1.9}

The distribution of heights of all students in a large university has a normal distribution with a mean of 66 in. and a standard deviation of 2 in. What is the probability that the mean height of 26 randomly selected students from this university will be more than 67 in.?


\section*{Solution:}

Let $X$ be the height. $X \sim N(\mu = 66, \sigma^2 = 2^2)$, and 
\[
\frac{\bar{X} - 66}{\frac{2}{\sqrt{26}}} = Z \sim N(0, 1^2).
\]
Then,
\[
P(\bar{X} > 170) = P\left(Z > 2\sqrt{26}\right) = 0.
\]


\section*{4.2.5}

Let $X_1, X_2, \dots, X_5$ be a random sample from the normal distribution with mean 55 and variance 223. Let
\[
Y = \frac{\sum_{i=1}^5 (X_i - 55)^2}{223}
\]
and
\[
Z = \frac{\sum_{i=1}^5 (X_i - \bar{X})^2}{223}.
\]

\begin{enumerate}
    \item[(a)] Find the distribution of the random variables $Y$ and $Z$.
    \item[(b)] Are $Y$ and $Z$ independent?
    \item[(c)] Find
    \begin{enumerate}
        \item[(i)] $P(0.554 < Y < 0.831)$, and
        \item[(ii)] $P(0.297 < Z < 0.484)$.
    \end{enumerate}
\end{enumerate}


\section*{Solution:}


\subsection*{(a) Distribution of the random variables $Y$ and $Z$}

\textbf{Distribution of $Y$:}  
The term $(X_i - 55)^2$ arises because $X_i \sim N(\mu = 55, \sigma^2 = 223)$.  
The scaled sum of squared deviations, $\sum_{i=1}^5 (X_i - 55)^2$, follows a chi-square distribution with degrees of freedom equal to the sample size ($n = 5$).  
Thus:
\[
Y = \frac{\sum_{i=1}^5 (X_i - 55)^2}{223} \sim \chi^2_5.
\]

\textbf{Distribution of $Z$:}  
The term $(X_i - \bar{X})^2$ represents deviations from the sample mean $\bar{X}$.  
The sum of squared deviations, $\sum_{i=1}^5 (X_i - \bar{X})^2$, follows a chi-square distribution with $n - 1$ degrees of freedom (since one parameter, the mean, is estimated).  
Thus:
\[
Z = \frac{\sum_{i=1}^5 (X_i - \bar{X})^2}{223} \sim \chi^2_4.
\]

\subsection*{(b) Are $Y$ and $Z$ independent?}

Since $Y$ involves deviations from the population mean (55), and $Z$ involves deviations from the sample mean ($\bar{X}$), their sources of variability are distinct.  
For a random sample from a normal distribution, the total variance can be partitioned into between-sample and within-sample variance.  
Thus:
\[
Y \text{ and } Z \text{ are independent.}
\]

\subsection*{(c) Calculations}

\begin{itemize}
    \item[(i)] $P(0.62 \leq Y \leq 0.76) = 0.0075$
    \item[(ii)] $P(0.77 \leq Z \leq 0.95) = 0.0251$
\end{itemize}


\section*{4.2.6}

Let $X$ and $Y$ be independent chi-square random variables with 14 and 5 degrees of freedom, respectively. Find:
\begin{enumerate}
    \item[(a)] $P(|X - Y| \leq 11.15)$,
    \item[(b)] $P(|X - Y| \geq 3.8)$.
\end{enumerate}

\section*{Solution:}


Let $X$ and $Y$ be independent chi-square random variables with 14 and 5 degrees of freedom, respectively. The random variable $D = X - Y$ can be approximated as normally distributed due to the central limit theorem.

\subsection*{(a) Find $P(|X - Y| \leq 11.15)$}

\textbf{Step 1: Mean and variance of $D$}

\[
\mu_X = 14, \quad \sigma_X^2 = 2 \times 14 = 28,
\]
\[
\mu_Y = 5, \quad \sigma_Y^2 = 2 \times 5 = 10.
\]

For $D = X - Y$:
\[
\mu_D = \mu_X - \mu_Y = 14 - 5 = 9,
\]
\[
\sigma_D^2 = \sigma_X^2 + \sigma_Y^2 = 28 + 10 = 38, \quad \sigma_D = \sqrt{38} \approx 6.16.
\]

Thus, $D \sim N(9, 38)$.

\textbf{Step 2: Standardize the inequality}

We need:
\[
P(|X - Y| \leq 11.15) = P(-11.15 \leq D \leq 11.15).
\]

Convert to the standard normal distribution ($Z$):
\[
Z = \frac{D - \mu_D}{\sigma_D}.
\]

Substitute:
\[
P(-11.15 \leq D \leq 11.15) = P\left(\frac{-11.15 - 9}{6.16} \leq Z \leq \frac{11.15 - 9}{6.16}\right).
\]

Simplify:
\[
P(-11.15 \leq D \leq 11.15) = P(-3.27 \leq Z \leq 0.35).
\]

\textbf{Step 3: Use the standard normal table}

From the standard normal table:
\[
P(Z \leq -3.27) \approx 0.0005, \quad P(Z \leq 0.35) \approx 0.6368.
\]

Thus:
\[
P(-3.27 \leq Z \leq 0.35) = P(Z \leq 0.35) - P(Z \leq -3.27),
\]
\[
P(-3.27 \leq Z \leq 0.35) = 0.6368 - 0.0005 = 0.6363.
\]

\textbf{Final Answer for (a):}
\[
P(|X - Y| \leq 11.15) \approx 0.6363.
\]

\subsection*{(b) Find $P(|X - Y| \geq 3.8)$}

\textbf{Step 1: Rewrite the inequality}

\[
P(|X - Y| \geq 3.8) = 1 - P(|X - Y| < 3.8).
\]

\textbf{Step 2: Standardize the inequality}

\[
P(|X - Y| < 3.8) = P(-3.8 \leq D \leq 3.8).
\]

Standardize:
\[
P(-3.8 \leq D \leq 3.8) = P\left(\frac{-3.8 - 9}{6.16} \leq Z \leq \frac{3.8 - 9}{6.16}\right).
\]

Simplify:
\[
P(-3.8 \leq D \leq 3.8) = P(-2.07 \leq Z \leq -0.84).
\]

\textbf{Step 3: Use the standard normal table}

From the standard normal table:
\[
P(Z \leq -2.07) \approx 0.0192, \quad P(Z \leq -0.84) \approx 0.2005.
\]

Thus:
\[
P(-2.07 \leq Z \leq -0.84) = P(Z \leq -0.84) - P(Z \leq -2.07),
\]
\[
P(-2.07 \leq Z \leq -0.84) = 0.2005 - 0.0192 = 0.1813.
\]

\textbf{Step 4: Find $P(|X - Y| \geq 3.8)$}

\[
P(|X - Y| \geq 3.8) = 1 - P(-3.8 \leq D \leq 3.8),
\]
\[
P(|X - Y| \geq 3.8) = 1 - 0.1813 = 0.8187.
\]

\textbf{Final Answer for (b):}
\[
P(|X - Y| \geq 3.8) \approx 0.8187.
\]


\section*{4.2.15}

Find the variance of $S^2$, assuming the sample $X_1, X_2, \dots, X_n$ is from $N(\mu, \sigma^2)$.

\section*{Solution:}

If $X_1, X_2, \dots, X_n$ is from $N(\mu, \sigma^2)$,  
then, by Theorem 4.2.8, $\frac{(n-1)S^2}{\sigma^2}$ is from $\chi^2_{(n-1)}$.

Then, by Exercise 4.2.13:
\[
\text{Var} \left[ \frac{(n-1)S^2}{\sigma^2} \right] = 2(n-1).
\]

Since $\text{Var}(aX) = a^2\text{Var}(X)$:
\[
\text{Var} \left( \frac{(n-1)^2}{\sigma^4} \text{Var}(S^2) \right) = 2(n-1).
\]

Simplifying after multiplying by $\frac{\sigma^4}{(n-1)^2}$, we obtain:
\[
\text{Var}(S^2) = \frac{2\sigma^4}{n-1}.
\]

\section*{4.2.23}

Let $X_1, X_2, \dots, X_{n_1}$ be a random sample from a normal population with mean $\mu_1$ and variance $\sigma^2$, and $Y_1, Y_2, \dots, Y_{n_2}$ be a random sample from an independent normal population with mean $\mu_2$ and variance $\sigma^2$. Show that
\[
T = \frac{\left(\bar{X} - \bar{Y}\right) - \left(\mu_1 - \mu_2\right)}{\sqrt{\frac{(n_1 - 1)S_1^2 + (n_2 - 1)S_2^2}{n_1 + n_2 - 2} \left(\frac{1}{n_1} + \frac{1}{n_2}\right)}} \sim T_{(n_1 + n_2 - 2)}.
\]

\section*{Solution:}

If $X_1, X_2, \dots, X_{n_1}$ is a random sample from a normal population with mean $\mu_1$ and variance $\sigma^2$, and if $Y_1, Y_2, \dots, Y_{n_2}$ is a random sample from an independent normal population with mean $\mu_2$ and variance $\sigma^2$, then:
\[
\bar{X} \sim N\left(\mu_1, \frac{\sigma^2}{n_1}\right), \quad \bar{Y} \sim N\left(\mu_2, \frac{\sigma^2}{n_2}\right), \quad \frac{(n_1 - 1)S_1^2}{\sigma^2} \sim \chi^2_{(n_1 - 1)}, \quad \frac{(n_2 - 1)S_2^2}{\sigma^2} \sim \chi^2_{(n_2 - 1)}.
\]

Then:
\[
\bar{X} - \bar{Y} \sim N\left(\mu_1 - \mu_2, \frac{\sigma^2}{n_1} + \frac{\sigma^2}{n_2}\right), \quad \frac{(n_1 - 1)S_1^2}{\sigma^2} + \frac{(n_2 - 1)S_2^2}{\sigma^2} \sim \chi^2_{(n_1 + n_2 - 2)}.
\]

Thus:
\[
\frac{\bar{X} - \bar{Y} - (\mu_1 - \mu_2)}{\sqrt{\frac{\sigma^2}{n_1} + \frac{\sigma^2}{n_2}}} \sim N(0, 1),
\]
and:
\[
\frac{(n_1 - 1)S_1^2}{\sigma^2} + \frac{(n_2 - 1)S_2^2}{\sigma^2} \sim \chi^2_{(n_1 + n_2 - 2)}.
\]

Then, since the samples are independent, we have by definition that:
\[
\frac{\bar{X} - \bar{Y} - (\mu_1 - \mu_2)}{\sqrt{\frac{\sigma^2}{n_1} + \frac{\sigma^2}{n_2}} \cdot \sqrt{\frac{(n_1 - 1)S_1^2}{\sigma^2} + \frac{(n_2 - 1)S_2^2}{\sigma^2}} \Big/ \sqrt{n_1 + n_2 - 2}} \sim T_{(n_1 + n_2 - 2)}.
\]

This simplifies to:
\[
\frac{\bar{X} - \bar{Y} - (\mu_1 - \mu_2)}{\sqrt{\frac{(n_1 - 1)S_1^2 + (n_2 - 1)S_2^2}{n_1 + n_2 - 2} \cdot \left(\frac{1}{n_1} + \frac{1}{n_2}\right)}} \sim T_{(n_1 + n_2 - 2)}. \quad \text{Q.E.D.}
\]


\section*{4.2.26}

Let the random variables $X_1, X_2, \dots, X_{10}$ be normally distributed with mean $8$ and variance $4$. Find a number $a$ such that
\[
P\left( \sum_{i=1}^{10} \left(\frac{X_i - 8}{2}\right)^2 \leq a \right) = 0.95.
\]

\section*{Solution:}

Let $X$ and $Y$ be independent chi-square random variables with 14 and 5 degrees of freedom, respectively. The random variable $D = X - Y$ can be approximated as normally distributed due to the central limit theorem.

\subsection*{(a) Find $P(|X - Y| \leq 11.15)$}

\textbf{Step 1: Mean and variance of $D$}

\[
\mu_X = 14, \quad \sigma_X^2 = 2 \times 14 = 28,
\]
\[
\mu_Y = 5, \quad \sigma_Y^2 = 2 \times 5 = 10.
\]

For $D = X - Y$:
\[
\mu_D = \mu_X - \mu_Y = 14 - 5 = 9,
\]
\[
\sigma_D^2 = \sigma_X^2 + \sigma_Y^2 = 28 + 10 = 38, \quad \sigma_D = \sqrt{38} \approx 6.16.
\]

Thus, $D \sim N(9, 38)$.

\textbf{Step 2: Standardize the inequality}

We need:
\[
P(|X - Y| \leq 11.15) = P(-11.15 \leq D \leq 11.15).
\]

Convert to the standard normal distribution ($Z$):
\[
Z = \frac{D - \mu_D}{\sigma_D}.
\]

Substitute:
\[
P(-11.15 \leq D \leq 11.15) = P\left(\frac{-11.15 - 9}{6.16} \leq Z \leq \frac{11.15 - 9}{6.16}\right).
\]

Simplify:
\[
P(-11.15 \leq D \leq 11.15) = P(-3.27 \leq Z \leq 0.35).
\]

\textbf{Step 3: Use the standard normal table}

From the standard normal table:
\[
P(Z \leq -3.27) \approx 0.0005, \quad P(Z \leq 0.35) \approx 0.6368.
\]

Thus:
\[
P(-3.27 \leq Z \leq 0.35) = P(Z \leq 0.35) - P(Z \leq -3.27),
\]
\[
P(-3.27 \leq Z \leq 0.35) = 0.6368 - 0.0005 = 0.6363.
\]

\textbf{Final Answer for (a):}
\[
P(|X - Y| \leq 11.15) \approx 0.6363.
\]

\subsection*{(b) Find $P(|X - Y| \geq 3.8)$}

\textbf{Step 1: Rewrite the inequality}

\[
P(|X - Y| \geq 3.8) = 1 - P(|X - Y| < 3.8).
\]

\textbf{Step 2: Standardize the inequality}

\[
P(|X - Y| < 3.8) = P(-3.8 \leq D \leq 3.8).
\]

Standardize:
\[
P(-3.8 \leq D \leq 3.8) = P\left(\frac{-3.8 - 9}{6.16} \leq Z \leq \frac{3.8 - 9}{6.16}\right).
\]

Simplify:
\[
P(-3.8 \leq D \leq 3.8) = P(-2.07 \leq Z \leq -0.84).
\]

\textbf{Step 3: Use the standard normal table}

From the standard normal table:
\[
P(Z \leq -2.07) \approx 0.0192, \quad P(Z \leq -0.84) \approx 0.2005.
\]

Thus:
\[
P(-2.07 \leq Z \leq -0.84) = P(Z \leq -0.84) - P(Z \leq -2.07),
\]
\[
P(-2.07 \leq Z \leq -0.84) = 0.2005 - 0.0192 = 0.1813.
\]

\textbf{Step 4: Find $P(|X - Y| \geq 3.8)$}

\[
P(|X - Y| \geq 3.8) = 1 - P(-3.8 \leq D \leq 3.8),
\]
\[
P(|X - Y| \geq 3.8) = 1 - 0.1813 = 0.8187.
\]

\textbf{Final Answer for (b):}
\[
P(|X - Y| \geq 3.8) \approx 0.8187.
\]


\section*{4.3.1}

The lifetime $X$ of a certain electrical fuse has the following pdf:
\[
f(x) = 
\begin{cases} 
\frac{1}{10}e^{-x/10}, & x > 0, \\
0, & \text{otherwise}.
\end{cases}
\]

Suppose two such fuses are in series and operate independently in a system. Find the pdf of the lifetime $Y$ of the system. (The system will work only if both of the fuses operate.)


\section*{Solution:}

The probability density function (pdf) of $X$ is:
\[
f(x) = \frac{1}{10}e^{-x/10}, \; x > 0.
\]

The cumulative distribution function (CDF) of $\tau_1$ is:
\[
F_{\tau_1}(t) = \int_0^t \frac{1}{10} e^{-x/10} \, dx = -e^{-x/10} \Big|_0^t = 1 - e^{-t/10}.
\]

Let $Y$ represent the life length of the system. Since the system operates as long as both fuses work, 
\[
Y = \min(\tau_1, \tau_2),
\]
where $\tau_1$ and $\tau_2$ are independent lifetimes of the two fuses. The CDF of $Y$ is:
\[
F_Y(y) = 1 - [1 - F_{\tau_1}(y)]^2.
\]

The pdf of $Y$ is given by:
\[
f_Y(y) = 2 f_{\tau_1}(y) [1 - F_{\tau_1}(y)],
\]
which simplifies to:
\[
f_Y(y) = 
\begin{cases} 
\frac{1}{5}e^{-y/5}, & 0 < y < \infty, \\
0, & \text{otherwise}.
\end{cases}
\]

\section*{4.3.4}

Let $X_1, \dots, X_{10}$ be a random sample from $U[0, 1]$. Find the joint density of $Y_2$ and $Y_7$, where $Y_i, i = 1, 2, \dots, 10$ are order statistics of $X_1, \dots, X_{10}$.


\section*{Solution:}

\subsection*{Step 1: Order Statistics}
The order statistics are defined as the sorted values of the sample:
\[
Y_1 \leq Y_2 \leq \dots \leq Y_{10},
\]
where $Y_2$ is the second smallest value and $Y_7$ is the seventh smallest value.

\subsection*{Step 2: Joint Distribution of Order Statistics}
The joint probability density function (PDF) of two order statistics $Y_2$ and $Y_7$ for a sample of size $n = 10$ is given by:
\[
f_{Y_2, Y_7}(y_2, y_7) = 
\begin{cases} 
\frac{n!}{(k-1)!(m-k-1)!(n-m)!} \, y_2^{k-1} (y_7 - y_2)^{m-k-1} (1 - y_7)^{n-m}, & 0 \leq y_2 \leq y_7 \leq 1, \\
0, & \text{otherwise}.
\end{cases}
\]
where:
\begin{itemize}
    \item $k = 2$ (index of $Y_2$),
    \item $m = 7$ (index of $Y_7$),
    \item $n = 10$ (sample size).
\end{itemize}

\subsection*{Step 3: Substitution and Simplification}
Substitute $n = 10$, $k = 2$, and $m = 7$ into the formula:
\[
f_{Y_2, Y_7}(y_2, y_7) = 
\begin{cases} 
\frac{10!}{(2-1)!(7-2-1)!(10-7)!} \, y_2^{1} (y_7 - y_2)^{4} (1 - y_7)^{3}, & 0 \leq y_2 \leq y_7 \leq 1, \\
0, & \text{otherwise}.
\end{cases}
\]

Simplify the factorial terms:
\[
\frac{10!}{1! \cdot 4! \cdot 3!} = \frac{362880}{1 \cdot 24 \cdot 6} = 2520.
\]

Thus:
\[
f_{Y_2, Y_7}(y_2, y_7) = 
\begin{cases} 
2520 \, y_2 (y_7 - y_2)^4 (1 - y_7)^3, & 0 \leq y_2 \leq y_7 \leq 1, \\
0, & \text{otherwise}.
\end{cases}
\]

\subsection*{Step 4: Final Answer}
The joint density of $Y_2$ and $Y_7$ is:
\[
f_{Y_2, Y_7}(y_2, y_7) = 
\begin{cases} 
2520 \, y_2 (y_7 - y_2)^4 (1 - y_7)^3, & 0 \leq y_2 \leq y_7 \leq 1, \\
0, & \text{otherwise}.
\end{cases}
\]

\section*{4.3.9}

Let $X_1, \dots, X_n$ be a random sample from a normal population with mean 10 and variance 4. What is the probability that the largest observation is greater than 10?


\section*{Solution:}

Let $X_1, \dots, X_n$ be a random sample from $N(10, 4)$. The probability that the largest observation is greater than 10 is:
\[
P(Y_n > 10) = 1 - P(Y_n \leq 10).
\]

The CDF $F_n(y)$ of $Y_n$ is $[F(y)]^n$, where $F(y)$ is the CDF of $X$ evaluated at $y$.

Thus:
\[
P(Y_n \leq y) = F_n(y) = [F(y)]^n = [P(X \leq y)]^n.
\]

For the normal distribution:
\[
P(X \leq y) = P\left(Z \leq \frac{y - 10}{2}\right).
\]

Substituting:
\[
P(Y_n \leq y) = \left[P\left(Z \leq \frac{y - 10}{2}\right)\right]^n.
\]

Thus:
\[
P(Y_n > y) = 1 - P(Y_n \leq y),
\]
\[
P(Y_n > y) = 1 - \left[P\left(Z \leq \frac{y - 10}{2}\right)\right]^n.
\]

For \( y = 10 \), this simplifies to:
\[
P(Y_n > 10) = 1 - \left[P(Z \leq 0)\right]^n,
\]
\[
P(Y_n > 10) = 1 - (0.5)^n.
\]

\newpage % Forces the following content to appear on a new page

\section*{\LARGE Projects}

\section*{Project 1:}

\begin{enumerate}
    \item Generating normal random samples (suppose $\mu$, $\sigma^2$ fixed) calculate the median, sample mean, and variance with different sample sizes, for example, $n = 100$, $500$, or $1000$ and compare with the mean and variance of the population. Draw a table or figure to show them, if necessary.
\end{enumerate}


\section*{Solution:}

\lstset{
    basicstyle=\ttfamily\footnotesize,
    keywordstyle=\color{blue},
    stringstyle=\color{red},
    commentstyle=\color{gray},
    showstringspaces=false,
    breaklines=true,
    frame=single,
    numbers=left,
    numberstyle=\tiny,
    captionpos=b,
    language=Python
}



\section*{Python Code Generating and Analyze Normal Samples}

\lstset{language=Python}
\begin{lstlisting}
import numpy as np
import pandas as pd
import matplotlib.pyplot as plt

# Parameters of the population
mu = 50  # Population mean
sigma2 = 25  # Population variance
sigma = np.sqrt(sigma2)  # Population standard deviation

# Sample sizes
sample_sizes = [100, 500, 1000]

# To store results
results = []

# Generate samples and calculate statistics
for n in sample_sizes:
    # Generate a normal random sample
    sample = np.random.normal(mu, sigma, n)

    # Calculate statistics
    sample_mean = np.mean(sample)
    sample_median = np.median(sample)
    sample_variance = np.var(sample, ddof=1)  # Sample variance (unbiased)

    # Append results
    results.append({
        "Sample Size": n,
        "Sample Mean": sample_mean,
        "Sample Median": sample_median,
        "Sample Variance": sample_variance
    })

# Convert results to a DataFrame
df_results = pd.DataFrame(results)

# Display the results in a table
print("Sample Statistics Comparison")
print(df_results)

# Plot the sample mean and variance compared to the population
plt.figure(figsize=(12, 6))

# Plot sample mean
plt.subplot(1, 2, 1)
plt.plot(sample_sizes, df_results["Sample Mean"], marker='o', label="Sample Mean")
plt.axhline(y=mu, color='r', linestyle='--', label="Population Mean")
plt.xlabel("Sample Size")
plt.ylabel("Mean")
plt.title("Sample Mean vs Population Mean")
plt.legend()

# Plot sample variance
plt.subplot(1, 2, 2)
plt.plot(sample_sizes, df_results["Sample Variance"], marker='o', label="Sample Variance")
plt.axhline(y=sigma2, color='r', linestyle='--', label="Population Variance")
plt.xlabel("Sample Size")
plt.ylabel("Variance")
plt.title("Sample Variance vs Population Variance")
plt.legend()

plt.tight_layout()
plt.show()
\end{lstlisting}


\begin{figure}[H] % The H option forces the figure to be placed exactly here
\centering

\includegraphics[width=1\textwidth]{D:/Masters NPU Semeter/3rd Semester/Statistics/Homework/Submit/Homework 4/figs/project_1a.png}
\caption{Project 1}
\end{figure}

\begin{figure}[H] % The H option forces the figure to be placed exactly here
\centering

\includegraphics[width=1\textwidth]{D:/Masters NPU Semeter/3rd Semester/Statistics/Homework/Submit/Homework 4/figs/project_1b.png}
\caption{Project 1}
\end{figure}

\section*{Project 2:}
\begin{enumerate}
    \item Drawing 10,000 samples, each sample of size 15 (suppose $\mu$, $\sigma^2$ fixed)\\
    Compute the mean of each sample and draw a histogram for the means. This will be an approximate sampling distribution of $\bar{X}$ for a fixed sample of size 15, and compare with the pdf figure of $\bar{X} \sim N(\mu, \sigma^2 / 15)$.
\end{enumerate}

\section*{Solution:}

\lstset{
    basicstyle=\ttfamily\footnotesize,
    keywordstyle=\color{blue},
    stringstyle=\color{red},
    commentstyle=\color{gray},
    showstringspaces=false,
    breaklines=true,
    frame=single,
    numbers=left,
    numberstyle=\tiny,
    captionpos=b,
    language=Python
}

\section*{Python Code for Sampling Distribution of the Sample Mean}

\lstset{language=Python}
\begin{lstlisting}
import numpy as np
import matplotlib.pyplot as plt
from scipy.stats import norm

# Parameters of the population
mu = 50  # Population mean
sigma2 = 25  # Population variance
sigma = np.sqrt(sigma2)  # Population standard deviation
sample_size = 15  # Sample size
num_samples = 10000  # Number of samples

# Generate 10,000 samples, each of size 15
sample_means = []
for _ in range(num_samples):
    sample = np.random.normal(mu, sigma, sample_size)
    sample_means.append(np.mean(sample))

# Convert sample means to a NumPy array
sample_means = np.array(sample_means)

# Compute the parameters of the sampling distribution of the sample mean
sampling_sigma = sigma / np.sqrt(sample_size)  # Standard deviation of the sampling distribution

# Plot histogram of sample means
plt.figure(figsize=(10, 6))
plt.hist(sample_means, bins=30, density=True, alpha=0.6, color='blue', label='Histogram of Sample Means')

# Plot theoretical normal distribution
x = np.linspace(mu - 4 * sampling_sigma, mu + 4 * sampling_sigma, 1000)
theoretical_pdf = norm.pdf(x, loc=mu, scale=sampling_sigma)
plt.plot(x, theoretical_pdf, 'r-', label='Theoretical PDF ($N(\\mu, \\sigma^2/15)$)')

# Add labels and legend
plt.title('Sampling Distribution of the Sample Mean (n=15)')
plt.xlabel('Sample Mean')
plt.ylabel('Density')
plt.legend()
plt.grid()
plt.show()
\end{lstlisting}

\begin{figure}[H] % The H option forces the figure to be placed exactly here
\centering

\includegraphics[width=1\textwidth]{D:/Masters NPU Semeter/3rd Semester/Statistics/Homework/Submit/Homework 4/figs/project_2.png}
\caption{Project 2}
\end{figure}

\end{document}